\section{Methods}
\label{s:methods}

\subsection{Setting, potential outcomes and estimands}

Units $i=1,\dots,N$ (counties) are partitioned into $G$ clusters (states); $g(i)$ denotes the
cluster of unit $i$ and the cluster sizes $n_g$ are bounded. Weeks are indexed by
$t=1-S,\dots,T$; weeks $1,\dots,T$ form the estimation window and the $S$ earlier weeks are
used only to build lags. $Y_{it}$ denotes reported new cases and $D_{it}\in\{0,1\}$ the policy,
with $D_{it}=0$ for $t\le 0$. The matrix $W=(W_{ij})$ is a known, non-negative, row-normalised
spatial weight matrix with zero diagonal, and $\bar D_{it}=\sum_j W_{ij}D_{jt}$ is the share of
neighbours under the policy. With known weekly generation-interval weights $w_1,\dots,w_S\ge 0$,
$\sum_s w_s=1$ \citep{wallinga2007generation}, we write
\begin{gather*}
L_{it}=\sum_{s=1}^{S} w_s Y_{i,t-s},\qquad
N_{it}=\sum_{s=1}^{S} w_s \sum_{j} W_{ij}Y_{j,t-s},\\
\Lambda_{it}(\phi)=L_{it}+\phi N_{it}+\omega_i ,
\end{gather*}
for the local, imported and total infection pressure; $\omega_i>0$ is a small known background
importation term. For $\theta=(\gamma_2,\dots,\gamma_T,\tau,\kappa,\phi)$, with $\gamma_1=0$, let
\begin{equation}
m_{it}(\theta)=\exp(\gamma_t+\tau D_{it}+\kappa\bar D_{it})\,\Lambda_{it}(\phi).
\label{eq:mean}
\end{equation}

Potential outcomes $Y_{it}(\mathbf d)$ are indexed by the whole policy array
$\mathbf d=(\mathbf d_1,\dots,\mathbf d_T)$, $\mathbf d_t\in\{0,1\}^N$, so that interference between
units is allowed \citep{hudgens2008interference}. $H_{t-1}(\mathbf d)$ is the potential history
of all units' outcomes up to week $t-1$, and $\Lambda_{it}(\mathbf d;\phi)$ is the pressure computed
from it. The unit effects $\alpha=(\alpha_1,\dots,\alpha_N)$ are unobserved. We work with the
filtration
\[
\F_t=\sigma\bigl(\alpha,\ Y_{js}\ (j\le N,\ s\le t),\ D_{js}\ (j\le N,\ s\le t+1)\bigr),
\]
which contains the policy for week $t+1$, decided at the end of week $t$. Instruments are
functions of the observable part of $\F_{t-1}$.

The target parameters are defined by the transmission model of Assumption A1 below. The
\emph{direct transmission effect} $\tau$ is the log ratio of a unit's conditional reproduction
number, $R_{it}(\mathbf d)=\E\{Y_{it}(\mathbf d)\mid\alpha,H_{t-1}(\mathbf d)\}/\Lambda_{it}(\mathbf d;\phi_0)$,
with and without its own policy, holding the history and the neighbours' policies fixed. The
\emph{policy spillover} $\kappa$ is the corresponding effect of the share of neighbours under the
policy, and the \emph{contagion coupling} $\phi$ is the relative infectiousness of neighbours'
recent cases. For a set $\mathcal P$ of post-adoption unit-weeks, the \emph{cases averted} are
\begin{equation}
\psi_{\mathcal P}=\E\sum_{(i,t)\in\mathcal P}\{Y_{it}(\mathbf 0)-Y_{it}\},
\label{eq:psi}
\end{equation}
where $\mathbf 0$ is the array with no policy anywhere; $\psi_{\mathcal P}$ includes propagation
through the renewal process and across units.

\subsection{Assumptions}

\begin{description}
\item[A1 (renewal model with parallel log-transmission trends).] For every policy array
$\mathbf d$, unit $i$ and week $t\ge 1$,
$\E\{Y_{it}(\mathbf d)\mid \alpha,H_{t-1}(\mathbf d)\}=\exp(\alpha_i+\gamma_{t}+\tau_0 d_{it}+\kappa_0\bar d_{it})\Lambda_{it}(\mathbf d;\phi_0)$.
\item[A2 (no anticipation and sequential ignorability).] (a) $Y_{it}(\mathbf d)$ depends on
$\mathbf d$ only through $\mathbf d_1,\dots,\mathbf d_t$. (b) For every $t$ and $\mathbf d$ with
$\Pr(D_{1:t}=\mathbf d_{1:t})>0$,
$\E\{Y_{it}(\mathbf d)\mid\alpha,H_{t-1}(\mathbf d),D_{1:t}=\mathbf d_{1:t}\}=\E\{Y_{it}(\mathbf d)\mid\alpha,H_{t-1}(\mathbf d)\}$.
\item[A3 (consistency).] $Y_{it}=Y_{it}(D)$.
\item[A4 (conditional independence across clusters).] For every $t$, conditionally on $\F_{t-1}$,
the cluster vectors $(Y_{it})_{i:g(i)=g}$, $g=1,\dots,G$, are independent. This rules out
infection that crosses a cluster boundary within the same week. Section~\ref{ss:weekwithin}
discusses what the assumption buys and what happens without it.
\item[A4$'$ (local dependence).] Let $\mathcal C$ be the graph on clusters in which $g$ and $h$
are adjacent if $W_{ij}>0$ for some $i\in g$, $j\in h$, or if the policy rule of $g$ uses data
from $h$; its maximal degree is bounded. (i) The cluster primitives (unit effects and pre-sample
outcomes, weekly outcome innovations, weekly policy randomisation) are independent across
clusters, and the week-$t$ outcome innovations are independent of all primitives dated before
week $t$ and of the week-$t$ policy randomisation. (ii) $D_{gt}$ is a
measurable function of the week-$t$ policy randomisation of $g$ and of the histories before week
$t$ of clusters within distance $\rho$ of $g$ in $\mathcal C$; $Y_{gt}$ is a measurable function of
the week-$t$ innovation and unit effects of $g$ and of the outcomes before week $t$ and the
week-$t$ policies of $g$ and of the clusters adjacent to $g$.
\item[A5 (positive pressure).] $\Lambda_{it}(\phi)\ge c>0$ for all $i$, $t$ and
$\phi\in[0,\bar\Phi]$, where $\bar\Phi$ bounds the parameter space for $\phi$ and is distinct from
the instrument reference value $\bar\phi$ defined in Section~\ref{ss:estimation}.
\item[A6 (regularity).] Web Appendix A.1: compact parameter space with $\theta_0$ interior;
moment dominance; existence of the limits $\bar g$, $\Gamma$, $\Omega$ defined below, with
$\Omega$ positive definite; $\bar g(\theta)=0$ only at $\theta_0$; $\Gamma$ of full rank.
\item[A7 (reporting).] Reported cases are $Y_{it}=c_ic_tI_{it}$, where $I_{it}$ are infections,
$c_i>0$ is a unit reporting factor and $c_t>0$ a week reporting factor.
\end{description}

Assumption A1 states that, in the absence of any policy, $\log R_{it}(\mathbf 0)=\alpha_i+\gamma_t$.
Untreated transmission paths are therefore parallel on the log scale, while levels may differ arbitrarily
and may be related to adoption. The conditional variance is unrestricted, so overdispersion and
clustering of cases are allowed. Assumption A2(b) is a mean version of sequential randomisation
\citep{robins2000msm}. Policies may depend on $\alpha$ and on the entire past of all units, which
is the feedback that motivates this paper, but not on the shocks of the week they apply to.
A4 allows arbitrary dependence within clusters and dependence across clusters through past
outcomes (contagion); it rules out common unmodelled shocks to neighbouring clusters within a
week beyond those absorbed by $\gamma_t$. A4$'$ implies A4 and is only used to verify the laws
of large numbers. It fails when policies respond to national aggregates, in which case the laws
of large numbers become high-level conditions. A5 holds by construction when $\omega_i\ge c$.

\begin{proposition}[Ascertainment]
\label{prop:ascertainment}
Suppose A1 holds for infections, with infection pressure $\Lambda^{I}_{it}$, and reported cases
satisfy A7. Suppose also that the reported pressure uses the background term
$\omega^{Y}_{it}=c_ic_{t-1}\omega_i{}^{I}$ and that either $\phi_0=0$ or $c_i=c_j$ whenever $W_{ij}>0$.
(a) If $c_{t-s}=c_{t-1}$ for $s=1,\dots,S$, then reported cases satisfy A1 with the same
$\alpha_i$, $\tau_0$, $\kappa_0$, $\phi_0$ and with $\gamma_t$ replaced by $\gamma_t+\log(c_t/c_{t-1})$.
(b) Otherwise, $|\log\Lambda^{Y}_{it}-\log(c_ic_{t-1}\Lambda^{I}_{it})|\le\max_{1\le s\le S}|\log(c_{t-s}/c_{t-1})|$,
so the error in the log reproduction number of week $t$ is bounded by the largest relative change
in ascertainment within the generation window.
\end{proposition}

Proposition~\ref{prop:ascertainment} is why time-varying testing intensity does not bias $\tau$
as long as it is common across units and smooth relative to the generation interval. When
reports are a binomial thinning of infections, the reported pressure is a noisy proxy for
$c_ic_{t-1}\Lambda^I_{it}$, and A1 holds for reported data only approximately; we treat this
case as a sensitivity analysis, not as a theorem.

\subsection{Quasi-differenced moments}
\label{ss:moments}

For $t=1,\dots,T-1$ let $r_{it}(\theta)=m_{it}(\theta)/m_{i,t+1}(\theta)$ and define the forward
quasi-difference \citep{chamberlain1992sequential,wooldridge1997multiplicative}
\begin{equation}
u_{it}(\theta)=Y_{i,t+1}\,r_{it}(\theta)-Y_{it}.
\label{eq:qd}
\end{equation}

\begin{lemma}[Observed-data model and moment]
\label{lem:moment}
Under A1--A3, $\E(Y_{it}\mid\F_{t-1})=\mu_{it}:=e^{\alpha_i}m_{it}(\theta_0)$ for $t\ge 1$, and
the quasi-difference \eqref{eq:qd} satisfies $\E\{u_{it}(\theta_0)\mid\F_{t-1}\}=0$ for $t=1,\dots,T-1$, provided
$\E|Y_{i,t+1}r_{it}(\theta_0)|<\infty$.
\end{lemma}

The unit effect $e^{\alpha_i}$ cancels because $r_{it}$ is a ratio of means, so $\alpha_i$ is
neither estimated nor needed. Let $z_{it}$ be a $q$-vector of instruments that are
$\F_{t-1}$-measurable. Cluster contributions and the sample moment are
\[
g_g(\theta)=\sum_{i:g(i)=g}\sum_{t=1}^{T-1} z_{it}u_{it}(\theta),\qquad
\hat g(\theta)=G^{-1}\sum_{g=1}^{G}g_g(\theta).
\]
Write $\varepsilon_{is}=Y_{is}-\mu_{is}$ and $r_{is}=r_{is}(\theta_0)$.

\begin{lemma}[Martingale representation]
\label{lem:mds}
Under A1--A3, $g_g(\theta_0)=\sum_{s=1}^{T}\xi_{gs}$ with
\begin{gather*}
\xi_{gs}=\sum_{i:g(i)=g}a_{is}\varepsilon_{is},\\
a_{is}=z_{i,s-1}r_{i,s-1}\ind(s\ge 2)-z_{is}\ind(s\le T-1),
\end{gather*}
where $a_{is}$ is $\F_{s-1}$-measurable.
Hence $\E(\xi_{gs}\mid\F_{s-1})=0$ and $\E(g_gg_g{}^{\top})=\sum_{s}\E(\xi_{gs}\xi_{gs}{}^{\top})$. If
also A4 holds, then $\E\{g_g(\theta_0)g_h(\theta_0)^{\top}\}=0$ for $g\ne h$.
\end{lemma}

Lemma~\ref{lem:mds} is the key to inference under lagged contagion. Neighbouring clusters can have
correlated outcomes because each one's cases seed the other's epidemic; this correlation is
visible in the numerical checks of Web Appendix B. Under A4, their moment contributions are nevertheless
uncorrelated, because each contribution is a sum of innovations weighted by predetermined
quantities. The ordinary cluster-robust variance is therefore valid under A4, without a spatial HAC
correction for $\hat\theta$; same-week cross-cluster transmission is excluded by that assumption.

\subsection{Identification}
\label{ss:ident}

Let $x_{it}(\phi)=(D_{it},\bar D_{it},N_{it}/\Lambda_{it}(\phi))^{\top}$ and $e_t$ the $t$th unit
vector of length $T-1$, and parametrise the week effects by $\delta_t=\gamma_{t+1}-\gamma_t$, so
that $\theta=(\delta^{\top},\beta^{\top})^{\top}$ with $\beta=(\tau,\kappa,\phi)^{\top}$ and
$p=T+2$. Then
\begin{equation}
\frac{\partial u_{it}(\theta)}{\partial\theta^{\top}}=Y_{i,t+1}r_{it}(\theta)\bigl\{-e_t{}^{\top},\,x_{it}(\phi)^{\top}-x_{i,t+1}(\phi)^{\top}\bigr\}.
\label{eq:jac}
\end{equation}
Here and below, $\lim$ denotes the limit as $G\to\infty$. Let $\bar g(\theta)=\lim G^{-1}\sum_g\E g_g(\theta)$ and $\Gamma=\partial\bar g(\theta_0)/\partial\theta^{\top}$.
With absorbing adoption, write $\Sw_{i,t+1}=\ind(D_{it}=0,D_{i,t+1}=1)$ for switchers.

\begin{theorem}[Identification]
\label{thm:ident}
Assume A1--A3, A5 and A6.
(a) $\bar g(\theta_0)=0$, and the Jacobian is
\begin{equation*}
\Gamma=\lim G^{-1}\sum_g\E\sum_{i\in g}\sum_{t}z_{it}\mu_{it}\{-e_t{}^{\top},x_{it}{}^{\top}-x_{i,t+1}{}^{\top}\}.
\end{equation*}
If $\operatorname{rank}\Gamma=p$, then $\theta_0$ is the unique zero of $\bar g$ in a
neighbourhood of $\theta_0$. If $\operatorname{rank}\partial\bar g(\theta)/\partial\theta^{\top}=r<p$
for all $\theta$ in a neighbourhood of $\theta_0$, then $\theta_0$ is not locally identified.
(b) Let adoption be absorbing, fix $(\kappa,\phi)=(\kappa_0,\phi_0)$, use the instruments
$s_{it}e_t$ and $s_{it}D_{it}$ with $s_{it}>0$ $\F_{t-1}$-measurable, and solve the week moments
for $\delta$. With $A_t=\lim G^{-1}\E\sum_i s_{it}D_{it}\mu_{it}$ and
$K_t=\lim G^{-1}\E\sum_i s_{it}\mu_{it}\Sw_{i,t+1}$, if $A_t>0$ and $K_t>0$ for some $t$, then
the profiled moment for $\tau$ is strictly increasing and $\tau_0$ is its unique root on
$\mathbb R$.
\end{theorem}

Part (a) is the moment-condition analogue of \citet{rothenberg1971identification}. Part (b) shows
where the information about $\tau$ comes from, which the next result makes explicit.

\begin{proposition}[Sources of identification]
\label{prop:sources}
Assume A1--A3 and absorbing adoption.
(a) If $D_{it}=D_{i,t+1}$, then $u_{it}(\theta)$ does not depend on $\tau$.
(b) The $\tau$-column of $\Gamma$ is
\begin{equation*}
-\lim G^{-1}\sum_g\E\sum_{i\in g}\sum_t z_{it}\mu_{it}\Sw_{i,t+1}.
\end{equation*}
(c) For the lead $D_{i,t+1}$, which is $\F_t$- but not $\F_{t-1}$-measurable, and any
$\F_{t-1}$-measurable $z_{it}$,
$\E\{z_{it}D_{i,t+1}u_{it}(\theta_0)\}=-\E\{z_{it}\cov(D_{i,t+1},\varepsilon_{it}\mid\F_{t-1})\}$.
\end{proposition}

Treated stayers carry no direct information on $\tau$; they enter only through the week effects,
which are pinned down by all units. The effect is learned from the growth of cases in switchers
relative to the common growth. Part (c) shows that the usual lead-placebo test is
\emph{invalid under feedback}. When adoption responds to high current cases,
$\cov(D_{i,t+1},\varepsilon_{it}\mid\F_{t-1})>0$ and the lead moment is negative even though
the model is correct. Leads are valid placebo instruments only when adoption dates are fixed at
least two weeks in advance. Under feedback we instead test the model with the over-identification
statistic of Corollary~\ref{cor:J}, using additional predetermined instruments.

\subsection{Estimation and inference}
\label{ss:estimation}

The default instruments are
$z_{it}=s_{it}\{e_t{}^{\top},D_{it},\bar D_{it},\bar\phi N_{it}/\Lambda_{it}(\bar\phi)\}^{\top}$
with scale weight $s_{it}=1/\{\Lambda_{it}(\bar\phi)+1\}$ and a reference value $\bar\phi$. This
gives $q=p$, and the estimator $\hat\theta$ solves $\hat g(\theta)=0$. The week moments solve in
closed form for given $\beta$:
\begin{equation}
e^{\hat\delta_t(\beta)}=\frac{\sum_i s_{it}Y_{i,t+1}q_{it}(\beta)}{\sum_i s_{it}Y_{it}},
\label{eq:profile}
\end{equation}
with $q_{it}(\beta)=e^{\tau(D_{it}-D_{i,t+1})+\kappa(\bar D_{it}-\bar D_{i,t+1})}\Lambda_{it}(\phi)/\Lambda_{i,t+1}(\phi)$,
so only a three-dimensional root search over $\beta$ is required. With
$\hat\Gamma=G^{-1}\sum_g\partial g_g(\hat\theta)/\partial\theta^{\top}$, computed analytically from
\eqref{eq:jac}, and $\hat\Omega=G^{-1}\sum_g g_g(\hat\theta)g_g(\hat\theta)^{\top}$, the variance
estimator is $\hat V=\{G/(G-1)\}\hat\Gamma^{-1}\hat\Omega\hat\Gamma^{-\top}/G$. Confidence intervals
use $t_{G-1}$ critical values, a small-cluster adjustment recommended by
\citet{cameron2015practitioner}. For over-identified versions with $q>p$, $\hat\theta$ minimises
$\hat g(\theta)^{\top}\hat W\hat g(\theta)$.

Efficient two-step GMM with all available instruments is not our primary estimator. With
$G\approx 49$ states and $T+2$ or more moments, the estimated optimal weight matrix is close to
singular and two-step standard errors are known to be too small \citep{windmeijer2005finite}.
Over-identification is used only for specification tests with a few additional instruments.

\begin{theorem}[Asymptotic properties]
\label{thm:asy}
Assume A1--A6, and either A4$'$ or the high-level laws of large numbers (L) of Web Appendix A.1.
Let $\hat W\to_p W$ positive definite (with $W=I$ when $q=p$).
(a) $\hat\theta\to_p\theta_0$.
(b) $\sqrt G(\hat\theta-\theta_0)=G^{-1/2}\sum_g\varphi_g+o_p(1)\to_d N(0,V)$ with influence function
$\varphi_g=-(\Gamma^{\top}W\Gamma)^{-1}\Gamma^{\top}Wg_g(\theta_0)$ and
$V=(\Gamma^{\top}W\Gamma)^{-1}\Gamma^{\top}W\Omega W\Gamma(\Gamma^{\top}W\Gamma)^{-1}$, where
\[
\Omega=\lim G^{-1}\sum_g\E\{g_g(\theta_0)g_g(\theta_0)^{\top}\}=\lim G^{-1}\sum_g\sum_s\E(\xi_{gs}\xi_{gs}{}^{\top}).
\]
In the just-identified case $\varphi_g=-\Gamma^{-1}g_g(\theta_0)$ and $V=\Gamma^{-1}\Omega\Gamma^{-\top}$.
(c) $\hat\Gamma\to_p\Gamma$ and $\hat\Omega\to_p\Omega$, so that $G\hat V\to_p V$.
\end{theorem}

\begin{corollary}[Over-identification test]
\label{cor:J}
Under the conditions of Theorem~\ref{thm:asy} with $q>p$ and $\hat W=\tilde\Omega^{-1}$ computed
at a consistent first-step estimator, $J=G\,\hat g(\hat\theta)^{\top}\hat W\hat g(\hat\theta)\to_d\chi^2_{q-p}$.
\end{corollary}

The central-limit step uses a martingale argument over the $T$ week blocks combined with a
conditional Lindeberg argument across clusters within each block (Lemma~\ref{lem:clt}, Web Appendix A). It
is self-contained, so we need not appeal to the nesting conditions of the general martingale
central limit theorem \citep{hall1980martingale}. With independent clusters the result reduces to
standard clustered-sample theory \citep{hansen2019clustered}. Instruments that depend on a
first-step estimate, such as $\bar\phi=\hat\phi$ from a first pass or a predicted adoption hazard,
do not change $V$ (Lemma~\ref{lem:generated}, Web Appendix A), because derivatives of predetermined instruments
are orthogonal to $u_{it}(\theta_0)$.

\begin{remark}[Working-optimal instruments]
\label{rem:optimal}
Ignoring the correlation between $u_{it}$ and $u_{i,t+1}$ and using a Poisson working variance,
$\E\{\partial u_{it}/\partial\theta\mid\F_{t-1}\}=\mu_{it}\{-e_t,x_{it}-\E(x_{i,t+1}\mid\F_{t-1})\}$
and $\operatorname{var}(u_{it}\mid\F_{t-1})=\mu_{it}\{1+\E(r_{it}\mid\F_{t-1})\}$, so the working-optimal
instrument $\{x_{it}-\E(x_{i,t+1}\mid\F_{t-1})\}/\{1+\E(r_{it}\mid\F_{t-1})\}$ is free of
$\alpha_i$. Its policy component is the predicted switching hazard, which can be estimated in a
first step (Lemma~5). Fully efficient instruments for sequential moment restrictions require
forward orthogonalisation \citep{chamberlain1992sequential,chamberlain2022feedback}, and we do not
derive the efficiency bound. Under a full Poisson likelihood more feedback-robust moments exist
\citep{bonhomme2025feedback}; under our conditional-mean model only quasi-differenced moments
are used.
\end{remark}

Four results are used in the simulation and application but are developed in Web Appendix A, where they can be read alongside their proofs. The comparators of Section~\ref{s:simulation} have pseudo-true values that differ from $\theta_0$: the two-step log-transmission regression accumulates a small-count term and a feedback term, and the Poisson fixed-effects estimator is inconsistent at fixed $T$ because its offset is the lagged outcome (Web Appendix A.9). With time-invariant heterogeneous effects the constant-effect estimand is a convex combination of the unit effects, so no unit receives negative weight; when effects build up after adoption the event-time specification of Theorem~\ref{thm:asy} recovers them bin by bin (Web Appendix A.10). Cases averted are defined by a counterfactual that propagates through the renewal equation and the network; the recursion leaves a unit's series unchanged until a policy reaches it, directly or through its neighbours, so no counterfactual start date has to be chosen, and inference targets the realised-path quantity whose expectation is the population target (Web Appendix A.11). Assumption A4 is about the innovations rather than the outcomes: neighbouring clusters share history through $\bar D$ and $\Lambda$, and Lemma~\ref{lem:mds} weights each moment contribution by a $\F_{s-1}$-measurable factor, so what A4 excludes is infection that crosses a cluster boundary inside the observation week. Section~\ref{ss:sameweek} measures the cost of that exclusion (Web Appendix A.12).

\subsection{Computation, variants and root selection}
\label{ss:computation}

The primary estimator, M0, uses the instruments of Section~\ref{ss:estimation} with the
pre-specified reference value $\bar\phi=0.5$. It is computed by a local root search for $\beta$
from a fixed starting value, with $\gamma$ profiled out by \eqref{eq:profile}. Three variants are
used for checks and are fixed in advance. M0-OI adds two predetermined instruments,
$s_{it}N_{i,t-1}/\Lambda_{i,t-1}(\bar\phi)$ and $s_{it}\log\{(Y_{i,t-1}+0.5)/(Y_{i,t-2}+0.5)\}$, and is
estimated by two-step GMM; its role is to provide the $J$ statistic of Corollary~\ref{cor:J}.
M0-ET replaces $\tau D_{it}$ by event-time effects $\tau_0,\dots,\tau_5$ and a pooled effect for
six or more weeks since adoption, as recommended in Remark~\ref{rem:dynamic}. M0$'$ imposes
$\kappa=\phi=0$ and corresponds to Theorem~\ref{thm:bias}(d).

Theorem~\ref{thm:ident} is a local result and A6 assumes that $\theta_0$ is interior. When
contagion between units is absent ($\phi_0=0$), the true value is on the boundary of the
parameter space. Estimating equations can then have several roots, of which at most one is
consistent \citep{small2000multiple}, and boundary values change the limiting distribution of
the estimator \citep{self1987boundary,andrews1999boundary}. We therefore pre-specified a
global version, M0-G. For each $\phi$ on a grid in $[0,3]$, $(\gamma,\tau,\kappa)$ are profiled
out and the remaining $\phi$-equation is evaluated. Every sign change on the grid is refined to
a root, and the boundary point $\phi=0$ is added as a candidate. Among the candidates, M0-G
selects the one with the smallest over-identified GMM criterion $Q$ computed with the M0-OI
instruments, following the idea of choosing among roots with a goodness-of-fit criterion
\citep{heyde1998multiple}. Inference is conditional on the selected candidate, and a boundary
selection is reported as such, without a Wald interval for $\phi$.

The simulation in Section~\ref{s:simulation} showed that this pre-specified rule is not
sufficient. We therefore also report a variant designed after interim results had been seen,
M0-GI, which is labelled post hoc wherever it appears. It uses the same candidates as M0-G,
but $Q$ is built from the M0 moments together with the scale-weighted $\phi$-instrument evaluated at
two other reference values, $s_{it}N_{it}/\Lambda_{it}(0.05)$ and $s_{it}N_{it}/\Lambda_{it}(2.0)$,
with $s_{it}=1/\{\Lambda_{it}(\bar\phi)+1\}$. The rationale
follows from Lemma~\ref{lem:moment}: moments formed with any predetermined instrument have mean
zero at $\theta_0$, whatever the reference value, whereas a spurious root that solves the
$\bar\phi$-instrument equation need not solve the equations at other reference values. Because
the selection rule was chosen after seeing results, its performance requires confirmation in
designs that were not used to choose it. The profiled $\phi$-moment over the grid, together with $Q$, shows whether more than one root
exists; both are reported for the application in Section~\ref{s:application}.

For applied work the rules amount to the following sequence, which we recommend reporting in
full. (1) Fit M0 from a fixed starting value and record the local solution. (2) Evaluate the
profiled $\phi$-equation on the implementation grid
\[
\begin{aligned}
\phi\in\{&0,0.02,0.05,0.1,0.15,0.2,0.25,0.3,0.4,0.5,0.6,\\
        &0.7,0.8,1,1.25,1.5,2,2.5,3\},
\end{aligned}
\]
refine every adjacent sign change, and include $\phi=0$ as a boundary candidate. (3) If the
profile has a single interior sign change, report M0 with its cluster-robust interval, together
with the profile as a diagnostic figure. (4) If there are several sign changes, report every
candidate root and the value of $Q$ at each, and report the rule's selection as conditional on that
choice; when the selected candidate is on the boundary, state that no interval for $\phi$ is
available. (5) Report the over-identification statistic $J$ in every case, because it has power
against some forms of misspecification but, as the simulation shows, little power against a spurious
root.
